\documentclass[14pt]{beamer}
\usepackage{pepslides}
\title{Free oscillations}
\lecturenumber{23}
\hypersetup{pdftitle={Lecture 23: Free oscillations}}

\definecolor{peplightblue}{RGB}{114,159,207}
\definecolor{pepdarkblue}{RGB}{52,101,164}

\begin{document}

\titleslide

\begin{frame}[t]
\heading{From last time\ldots}
\centering
\includegraphics[height=0.62\paperheight]{L10F2.pdf}
\credit{Dziewonski \& Anderson (1981)}
\end{frame}

\sectionslide{Operator formalism}

\begin{frame}
\eqscale{17}
\textbf{Linearised equations with self-gravitation and rotation:}
\begin{eqn}
\begin{gathered}
\rho\left(\frac{\partial^{2}u_{i}}{\partial t^{2}}+2\epsilon_{ijk}\Omega_{j}\frac{\partial u_{k}}{\partial t}
+\epsilon_{ijk}\epsilon_{klm}\Omega_{j}\Omega_{l}u_{m}\right)-\frac{\partial}{\partial x_{j}}\left(A_{ijkl}\frac{\partial u_{k}}{\partial x_{l}}\right)
-\rho\gamma_{i}^{1}=-\frac{\partial\overline{S}_{ij}}{\partial x_{j}},\\[\eqgap]
\gamma_{i}^{1}=-G\int_{M}\rho'\left(\frac{\delta_{ij}}{[(x_{k}-x_{k}')(x_{k}-x_{k}')]^{3/2}}-\frac{3(x_{i}-x_{i}')(x_{j}-x_{j}')}{[(x_{k}-x_{k}')(x_{k}-x_{k}')]^{5/2}}\right)(u_{j}-u_{j}')\dd^{3}\mathbf{x}',\\[\eqgap]
A_{ijkl}\frac{\partial u_{k}}{\partial x_{l}}\hat{n}_{j}=\overline{S}_{ij}\hat{n}_{j}.
\end{gathered}
\end{eqn}
\textbf{Simplifying notations will help a lot!}
\end{frame}

\begin{frame}
\eqscale{18}
For an arbitrary \textbf{test function}, we have:
\begin{eqn}
\begin{aligned}
\int_{M}\Bigg[&\rho\left(\frac{\partial^{2}u_{i}}{\partial t^{2}}+2\epsilon_{ijk}\Omega_{j}\frac{\partial u_{k}}{\partial t}+\epsilon_{ijk}\epsilon_{klm}\Omega_{j}\Omega_{l}u_{m}\right)\\
&-\frac{\partial}{\partial x_{j}}\left(A_{ijkl}\frac{\partial u_{k}}{\partial x_{l}}-\overline{S}_{ij}\right)-\rho\gamma_{i}^{1}\Bigg]w_{i}^{*}\dd^{3}\mathbf{x}=0.
\end{aligned}
\end{eqn}
Integrating by parts and using the boundary conditions:
\begin{eqn}
\begin{aligned}
\int_{M}\Bigg[&\rho\left(\frac{\partial^{2}u_{i}}{\partial t^{2}}+2\epsilon_{ijk}\Omega_{j}\frac{\partial u_{k}}{\partial t}+\epsilon_{ijk}\epsilon_{klm}\Omega_{j}\Omega_{l}u_{m}-\gamma_{i}^{1}\right)w_{i}^{*}\\
&+\left(A_{ijkl}\frac{\partial u_{k}}{\partial x_{l}}-\overline{S}_{ij}\right)\frac{\partial w_{i}^{*}}{\partial x_{j}}\Bigg]\dd^{3}\mathbf{x}=0.
\end{aligned}
\end{eqn}
If the latter identity holds for all test functions, then we can recover both the equations of motion and the boundary
conditions. This is the \textbf{weak formulation} of the problem.

Weak and strong formulations are equivalent up to technicalities.
\end{frame}

\begin{frame}
We define the \textbf{kinetic energy form}:
\begin{eqn}
\langle \mathbf{w}|P|\mathbf{u}\rangle=\int_{M}\rho w_{i}^{*}u_{i}\dd^{3}\mathbf{x}.
\end{eqn}
The acceleration term in the weak form can then be written:
\begin{eqn}
\langle \mathbf{w}|P|\ddot{\mathbf{u}}\rangle.
\end{eqn}
This \textbf{sesquilinear form} is associated with the linear operator:
\begin{eqn}
(P\mathbf{u})_{i}=\rho u_{i}.
\end{eqn}
\end{frame}

\begin{frame}
We define the \textbf{Coriolis form}:
\begin{eqn}
\langle \mathbf{w}|W|\mathbf{u}\rangle=\int_{M}2\rho\epsilon_{ijk}w_{i}^{*}\Omega_{j}u_{k}\dd^{3}\mathbf{x},
\end{eqn}
associated with the linear operator:
\begin{eqn}
(W\mathbf{u})_{i}=2\rho\epsilon_{ijk}\Omega_{j}u_{k}.
\end{eqn}
\end{frame}

\begin{frame}
\eqscale{20}
We define the \textbf{potential energy form}:
\begin{eqn}
\langle \mathbf{w}|H|\mathbf{u}\rangle=\int_{M}\left(\rho\epsilon_{ijk}\epsilon_{klm}w_{i}^{*}\Omega_{j}\Omega_{l}u_{m}+A_{ijkl}\frac{\partial w_{i}^{*}}{\partial x_{j}}\frac{\partial u_{k}}{\partial x_{l}}-\rho w_{i}^{*}\gamma_{i}^{1}\right)\dd^{3}\mathbf{x},
\end{eqn}
along with the linear form linked to the stress glut:
\begin{eqn}
\langle \mathbf{w}|\mathbf{f}\rangle=\int_{M}\overline{S}_{ij}\frac{\partial w_{i}^{*}}{\partial x_{j}}\dd^{3}\mathbf{x}.
\end{eqn}
These two terms incorporate \textbf{both} the equations of motion and the boundary conditions.
\end{frame}

\begin{frame}
The weak form of the equations of motion can now be written:
\begin{eqn}
\langle \mathbf{w}|P|\ddot{\mathbf{u}}\rangle+\langle \mathbf{w}|W|\dot{\mathbf{u}}\rangle+\langle \mathbf{w}|H|\mathbf{u}\rangle=\langle \mathbf{w}|\mathbf{f}\rangle,
\end{eqn}
which is to hold for all test functions, $\mathbf{w}$.

We can formally think of the problem as
\begin{eqn}
P\ddot{\mathbf{u}}+W\dot{\mathbf{u}}+H\mathbf{u}=\mathbf{f},
\end{eqn}
but the boundary conditions are then left implicit.
\end{frame}

\begin{frame}
The kinetic energy form is \textbf{self-adjoint}:
\begin{eqn}
\langle \mathbf{w}|P|\mathbf{u}\rangle=\langle \mathbf{u}|P|\mathbf{w}\rangle^{*}.
\end{eqn}
It is also positive for non-zero $\mathbf{u}$, and so defines an \textbf{inner product}.

The Coriolis form is \textbf{anti-self-adjoint}:
\begin{eqn}
\begin{aligned}
\langle \mathbf{w}|W|\mathbf{u}\rangle&=\int_{M}2\rho\epsilon_{ijk}w_{i}^{*}\Omega_{j}u_{k}\dd^{3}\mathbf{x}\\[\eqgap]
&=-\left(\int_{M}2\rho\epsilon_{ijk}u_{i}^{*}\Omega_{j}w_{k}\dd^{3}\mathbf{x}\right)^{*}=-\langle \mathbf{u}|W|\mathbf{w}\rangle^{*}.
\end{aligned}
\end{eqn}
\end{frame}

\begin{frame}
The potential energy form is \textbf{self-adjoint}:
\begin{eqn}
\langle \mathbf{w}|H|\mathbf{u}\rangle=\langle \mathbf{u}|H|\mathbf{w}\rangle^{*}.
\end{eqn}
Elastic term is clear due to the hyperelastic symmetry, $A_{ijkl}=A_{klij}$.

Centrifugal term follows by a simple calculation -- details in the handout.

Gravitational term takes more work but proceeds in the same manner -- details are not in the handout
and are not examinable.
\end{frame}

\sectionslide{Non-rotating eigenvalue problem}

\begin{frame}
For long-period seismology, rotation is a relatively small term that is often handled through perturbation theory.

Weak form in a non-rotating earth model:
\begin{eqn}
\langle \mathbf{w}|P|\ddot{\mathbf{u}}\rangle+\langle \mathbf{w}|H|\mathbf{u}\rangle=\langle \mathbf{w}|\mathbf{f}\rangle.
\end{eqn}
\textbf{Free oscillations} are time-harmonic solutions with no force term:
\begin{eqn}
\begin{gathered}
u_{i}(\mathbf{x},t)=s_{i}(\mathbf{x})\ee^{\ii\omega t},\\[\eqgap]
-\omega^{2}\langle \mathbf{w}|P|\mathbf{s}\rangle+\langle \mathbf{w}|H|\mathbf{s}\rangle=0.
\end{gathered}
\end{eqn}
This is an eigenvalue problem, with the squared eigenfrequency as the eigenvalue.
\end{frame}

\begin{frame}
You need to know how to show that:
\begin{enumerate}
\item Squared eigenfrequencies are real
\item Eigenfunctions in distinct eigenspaces are orthogonal:
\end{enumerate}
\begin{eqn}
\langle \mathbf{s}_{1}|P|\mathbf{s}_{2}\rangle=0.
\end{eqn}
Further properties to know but without proof:
\begin{enumerate}
\item Eigenspaces are all finite-dimensional
\item Countably many squared eigenfrequencies, bounded from below and accumulating only at infinity
\item Eigenfunctions form a \textbf{complete basis} for vector fields
\end{enumerate}
\end{frame}

\begin{frame}
We use $k$ to label the eigenspaces, and $m$ the eigenfunctions in each eigenspace.

Within each eigenspace, we can choose an orthonormal basis, and hence:
\begin{eqn}
\langle km|P|k'm'\rangle=\delta_{kk'}\delta_{mm'}.
\end{eqn}
Completeness of the eigenfunctions means that we can write:
\begin{eqn}
\mathbf{u}=\sum_{km}\langle km|P|\mathbf{u}\rangle|km\rangle.
\end{eqn}
This is the \textbf{eigenfunction expansion} of the vector field.
\end{frame}

\begin{frame}
Equations of motion in the frequency domain:
\begin{eqn}
-\omega^{2}\langle \mathbf{w}|P|\tilde{\mathbf{u}}\rangle+\langle \mathbf{w}|H|\tilde{\mathbf{u}}\rangle=\langle \mathbf{w}|\tilde{\mathbf{f}}\rangle.
\end{eqn}
Expanding the solution in the eigenfunction basis:
\begin{eqn}
\tilde{\mathbf{u}}=\sum_{k'm'}\langle k'm'|P|\tilde{\mathbf{u}}\rangle|k'm'\rangle.
\end{eqn}
Frequency-domain solution readily derived:
\begin{eqn}
\tilde{\mathbf{u}}=\sum_{km}\frac{\langle km|\tilde{\mathbf{f}}\rangle}{\omega_{k}^{2}-\omega^{2}}|km\rangle.
\end{eqn}
Time-domain solution follows by contour integration -- not examinable, but you will use it in the second problem set.
\end{frame}

\begin{frame}
\centering
\includegraphics[height=0.66\paperheight]{L11F1.pdf}

\raggedright
Ten-day record following a large earthquake (Tohoku 2011, recorded in the Black Forest, Germany).

The effects of attenuation are clearly seen within real data.
\end{frame}

\begin{frame}
\centering
\includegraphics[height=0.7\paperheight]{L11F2.pdf}

\raggedright
\textbf{Resonance peaks} (or poles) associated with each eigenfrequency; vertical lines mark those of PREM.

The eigenfrequencies depend only on Earth structure, avoiding the source--structure trade-off.
\end{frame}

\begin{frame}
A real record has a finite length, $T$:
\begin{eqn}
\mathbf{u}_{\mathrm{obs}}(t)=h(t)\mathbf{u}(t).
\end{eqn}
Its spectrum is a smoothed version of the true one:
\begin{eqn}
\tilde{\mathbf{u}}_{\mathrm{obs}}(\omega)\propto\int_{-\infty}^{\infty}\tilde{\mathbf{u}}(\omega')\tilde{h}(\omega-\omega')\dd\omega'.
\end{eqn}
The resonance peaks are broadened over a width of order $T^{-1}$.

For good resolution in the frequency domain, we need \textbf{long time series} -- and hence the largest earthquakes.
\end{frame}

\sectionslide{Gravitational stability}

\begin{frame}
The non-rotating eigenvalue problem:
\begin{eqn}
-\omega_{k}^{2}\langle \mathbf{w}|P|km\rangle+\langle \mathbf{w}|H|km\rangle=0.
\end{eqn}
Taking the test function to equal the eigenfunction:
\begin{eqn}
\omega_{k}^{2}=\langle km|H|km\rangle.
\end{eqn}
Sign of the squared eigenfrequency determined by the potential energy form.

\textbf{Positive values} imply \textbf{stable} oscillations about equilibrium.

\textbf{Negative values} imply \textbf{unstable}, exponentially growing perturbations.

\textbf{Zero values} correspond to \textbf{linear growth in time} (double pole at the origin).
\end{frame}

\begin{frame}
\eqscale{21}
In the non-gravitating case, the potential energy form simplifies such that:
\begin{eqn}
\langle \mathbf{u}|H|\mathbf{u}\rangle=\frac{1}{4}\int_{M}A_{ijkl}\left(\frac{\partial u_{i}}{\partial x_{j}}
+\frac{\partial u_{j}}{\partial x_{i}}\right)^{*}\left(\frac{\partial u_{k}}{\partial x_{l}}+\frac{\partial u_{l}}{\partial x_{k}}\right)\dd^{3}\mathbf{x}.
\end{eqn}
This vanishes for linearised rigid body motions, with $B_{ij}$ anti-symmetric:
\begin{eqn}
u_{i}(\mathbf{x})=a_{i}+B_{ij}x_{j}.
\end{eqn}
Six-dimensional eigenspace at zero frequency, but not excited by earthquakes:
\begin{eqn}
\langle \mathbf{w}|\mathbf{f}\rangle=\int_{M}\overline{S}_{ij}\frac{\partial w_{i}^{*}}{\partial x_{j}}\dd^{3}\mathbf{x}.
\end{eqn}
\end{frame}

\begin{frame}
\eqscale{21}
In the non-gravitating case, the potential energy form simplifies such that:
\begin{eqn}
\langle \mathbf{u}|H|\mathbf{u}\rangle=\frac{1}{4}\int_{M}A_{ijkl}\left(\frac{\partial u_{i}}{\partial x_{j}}
+\frac{\partial u_{j}}{\partial x_{i}}\right)^{*}\left(\frac{\partial u_{k}}{\partial x_{l}}+\frac{\partial u_{l}}{\partial x_{k}}\right)\dd^{3}\mathbf{x}.
\end{eqn}
Non-negative values so long as we have \textbf{uniform pointwise stability}:
\begin{eqn}
A_{ijkl}(\mathbf{x})e_{ij}^{*}e_{kl}\ge c_{0}\,e_{ij}^{*}e_{ij},\quad c_{0}>0.
\end{eqn}
Implies elastic wave speeds are all positive.

In the isotropic case, it also implies that the P-wave speed is greater than the S-wave speed.
\end{frame}

\begin{frame}
\begin{columns}[c]
\begin{column}{0.4\textwidth}
\centering
\begin{tikzpicture}
\fill[peplightblue] (0,0) circle (3.3cm);
\draw[pepdarkblue!70!black,thin] (0,0) circle (3.3cm);
\fill[pepdarkblue] (0,0) circle (2.6cm);
\draw[pepdarkblue!70!black,thin] (0,0) circle (2.6cm);
\foreach \a in {0,90,180,270}
  \draw[-{Latex[length=2.5mm,width=2mm]},semithick] (\a:3.8cm) -- (\a:2.62cm);
\end{tikzpicture}
\end{column}
\begin{column}{0.56\textwidth}
\setlength{\parskip}{0.8em}%
Quantitative treatment possible, but difficult; a simple case is in the second problem set.

Stability depends on the interplay between gravitational and elastic forces.

Elastic forces always resist further deformation.

Self-gravitation is \textbf{inherently destabilising}.

A spherically symmetric compression would lower the (always negative) gravitational binding energy, while
raising the elastic potential energy.

For realistic earth models such as PREM, calculations find no negative squared eigenfrequencies.
\end{column}
\end{columns}
\end{frame}

\begin{frame}[t]
\heading{What you need to know and be able to do}
\begin{enumerate}
\item[(i)] The weak formulation of the equations of motion using bra-ket-like notation, and the properties of the
different sesquilinear forms.
\item[(ii)] How to derive the reality of the squared eigenfrequencies and orthogonality of eigenfunctions in distinct
eigenspaces. State without proof the remaining important properties of the eigenvalue problem.
\item[(iii)] How to obtain a solution of the forced equations of motion in a non-rotating earth model as an
eigenfunction expansion in the frequency domain. You do not need to know how to recover the time-domain solution, but
may be asked to verify that such a solution works.
\item[(iv)] That eigenfrequencies are associated with resonance peaks in the amplitude spectra obtained from
seismograms, and why it is necessary to look at data from large earthquakes to obtain adequate resolution.
\item[(v)] What it means for a non-rotating earth model to be stable in terms of the eigenfrequencies. You should
understand the quantitative argument for the non-gravitating case, and the qualitative one when gravity is included.
\end{enumerate}
\end{frame}

\end{document}
